\section{Screening for every coupling}\label{sec:8.4}

In GaAs the dielectric function vanishes at \SI{36.35}{meV}, the energy of the longitudinal optical phonon (\cref{sec:5.5}). A rate written as the unscreened rate divided by $|\varepsilon|^2$, the form \cref{sec:4.4} derived for a probe of the electron density, would be infinite there. For a scalar that couples to nucleons it is finite, and the phonon it excites sits at \SI{36.35}{meV} rather than at the \SI{33.62}{meV} of the transverse mode, because the Born charges shift the pole and leave the coupling alone (\cref{we:8-nucleon}). The rule $1/|\varepsilon|^2$, and the loss function behind it, belong to couplings that see the charge. A dark photon sees the total charge of electrons and nuclei, an electrophilic mediator the electrons alone, a hadrophilic scalar the nucleons, and a spin coupling no charge at all, and each is screened differently.

\subsection{The proper response of a channel}\label{sec:8.4-proper}

Electrons and nuclei interact through one thing, the Coulomb energy $H_{\mathrm{C}}$ of \cref{eq:2.3-HC}, which couples them only through the charge-number density $n_Z = -n_e + \sum_\kappa Z_\kappa n_\kappa$ of \cref{sec:2.3}. Phonons, electron--phonon vertices and screening are consequences of it. A dark-matter particle that transfers $\vect{q}$ couples to the rows of \cref{eq:8.1-rows}, and in a crystal the rows it excites induce charge at $\vect{q} + \vect{G}$ for every reciprocal-lattice vector (\cref{sec:4.4}), with $\vect{q}$ in the first zone and the dark-matter vertex at the row $\vect{G}_q$ of \cref{eq:A.3-unreduced}. The set of momenta $\{\vect{q} + \vect{G}\}$ is the \emph{channel} of the external field. The Coulomb lines that carry a momentum of the channel are the ones that screen the dark-matter vertex. All others, at $\vect{q}' \ne \vect{q} + \vect{G}$, belong to the material: they bind the electrons, make the phonons and dress the electron--phonon vertex.

Follow the argument of \cref{sec:4.4} with every species at once. Let the external field couple to row $b$ with amplitude $\phi_b$, and write $\delta O_a = \delta\langle O_a(\vect{q})^\dagger\rangle/\Vol$ for the induced change of row $a$, so that the Kubo formula \cref{eq:4.1-kubo} reads $\delta O_a = R_{ab}\phi_b$. The induced charge-number density at $\vect{q} + \vect{G}$, $\delta\nd_{Z,\vect{G}}$, produces the Hartree potential energy $v_{\mathrm{C}}(|\vect{q}+\vect{G}|)\,\delta\nd_{Z,\vect{G}}$ on a unit charge, and on the target it acts as an extra field on the row $n_Z(\vect{q}+\vect{G})$. With electrons alone, $n_Z = -n_e$ and the two signs cancel, which is \cref{eq:4.4-phitot}. The \emph{proper response of the channel}, $\tilde R$, is the response to the external and Hartree fields together,
\begin{equation}
  \delta O_a = \sum_b\tilde R_{ab}\,\phi_b + \sum_{\vect{G}}\tilde R_{aZ_{\vect{G}}}\,v_{\mathrm{C}}(|\vect{q}+\vect{G}|)\,\delta\nd_{Z,\vect{G}} ,
  \label{eq:8.4-proper}
\end{equation}
where the subscript $Z_{\vect{G}}$ stands for the row $n_Z(\vect{q}+\vect{G})$, so that $\tilde R_{aZ_{\vect{G}}} = -\tilde R_{a e_{\vect{G}}} + \sum_\kappa Z_\kappa\tilde R_{a\kappa_{\vect{G}}}$. Electrons and nuclei enter on the same footing, as two kinds of charge. In the language of \cref{sec:4.8}, $\tilde R$ is the sum of the diagrams that cannot be cut by removing a single Coulomb line of the channel. By the argument of \cref{eq:4.8-removal} it is also the response of the Hamiltonian with the Coulomb terms of the channel deleted, which has the same equilibrium state in the thermodynamic limit. Every exact property of \cref{chap:4} therefore holds for $\tilde R$, the Lehmann form and positivity included.

Taking $a = Z_{\vect{G}}$ in \cref{eq:8.4-proper} closes the equations for the induced charge, and the result for any pair of rows is (\cref{app:D.27})
\begin{keyresult}[Screened response for any coupling]
\begin{equation}
  \begin{gathered}
  R_{ab}(\vect{q},\omega) = \tilde R_{ab} + \sum_{\vect{G},\vect{G}'}\tilde R_{aZ_{\vect{G}}}\,\bigl(\varepsilon^{-1}\bigr)_{\vect{G}\vect{G}'}\,v_{\mathrm{C}}(|\vect{q}+\vect{G}'|)\,\tilde R_{Z_{\vect{G}'}b} , \\
  \varepsilon_{\vect{G}\vect{G}'} = \delta_{\vect{G}\vect{G}'} - v_{\mathrm{C}}(|\vect{q}+\vect{G}|)\,\tilde R_{Z_{\vect{G}}Z_{\vect{G}'}} ,
  \end{gathered}
  \label{eq:8.4-screened}
\end{equation}
with every function at $(\vect{q},\omega)$. The matrix $\varepsilon$ is the total dielectric matrix of electrons and nuclei. A row is screened through its proper correlation with the charge, $\tilde R_{aZ}$, and through nothing else.
\end{keyresult}
\noindent Without local fields the sums reduce to $\vect{G} = 0$, and
\begin{equation}
  R_{ab} = \tilde R_{ab} + v_{\mathrm{C}}(q)\,\frac{\tilde R_{aZ}\,\tilde R_{Zb}}{\varepsilon} , \qquad \varepsilon = 1 - v_{\mathrm{C}}(q)\,\tilde R_{ZZ} .
  \label{eq:8.4-rank1}
\end{equation}
The reading is the one of \cref{sec:4.4}, with the vertex made explicit. The dark matter acts on row $b$; the target answers with a proper charge $\tilde R_{Zb}$; its field, screened by the whole medium, $\varepsilon^{-1}v_{\mathrm{C}}$, acts back on row $a$ through $\tilde R_{aZ}$. The dielectric function contains the polarization of the nuclei as well as that of the electrons, the static $\varepsilon_0$ of \cref{sec:5.5} at low frequency and $\varepsilon_\infty$ between the phonons and the gap. With the nuclei clamped it reduces to the electronic $\varepsilon$ of \cref{sec:4.4}. For time-dependent density-functional theory the same algebra holds with $f_{\mathrm{xc}}$ inside the proper part (\cref{sec:4.6}).

The proper cross responses are material quantities. The charge--charge block is the dielectric matrix that codes compute in the random-phase approximation with local fields or with an exchange--correlation kernel (\cref{app:E.6}), and its ionic part at long wavelength is fixed by the Born charges and phonons of density-functional perturbation theory (\cref{app:E.4}). The blocks $\tilde R_{aZ}$ for a current, spin or kinetic-energy row are computed from the same Kohn--Sham states as the response of row $a$ to a scalar potential.

\subsection{Four couplings}\label{sec:8.4-couplings}

The rows of a density coupling are the densities of the species, weighted by the coefficients $c_1^{(s)}$ of \cref{sec:8.1}: the dark matter acts through $\sum_sc_1^{(s)}n_s(\vect{q})$, with $c_1^{(\kappa)} = [Z_\kappa c_1^{(p)} + (A_\kappa - Z_\kappa)c_1^{(n)}]F_N$ from \cref{eq:8.1-nuclear}. Three choices of the coefficients cover the standard density couplings, a fourth coupling, to the spin, carries no charge at all, and \cref{eq:8.4-screened} treats each in one line.

\emph{Total charge.} A dark photon, heavy or light, couples to $-\tilde V n_Z$ (\cref{we:8-darkphoton}), and so does any millicharged particle. With $a = b = Z$, \cref{eq:8.4-screened} collapses to $v_{\mathrm{C}}R_{ZZ} = \varepsilon^{-1} - 1$ (\cref{app:D.27}), the relation \cref{sec:4.4} stated, now with the nuclei moving. The fluctuation--dissipation theorem then gives $S_{ZZ} = (2q^2/e^2)\,L/(1 - \ee^{-\beta\omega})$, with $L = \Imag[-(\varepsilon^{-1})_{00}]$ the loss function of the total dielectric matrix, exactly and at every energy. Below the gap $L$ contains the longitudinal optical phonons of a polar crystal, which \cref{sec:11.2} turns into a rate.

\emph{Electron density.} A scalar or vector mediator that couples to electrons and not to nucleons acts through $c_1^{(e)}n_e$, and \cref{eq:8.4-screened} gives
\begin{equation}
  R_{ee} = \tilde R_{ee} + \tilde R_{eZ}\,\varepsilon^{-1}v_{\mathrm{C}}\,\tilde R_{Ze} , \qquad
  \tilde R_{eZ} = -\tilde R_{ee} + \sum_\kappa Z_\kappa\tilde R_{e\kappa} .
  \label{eq:8.4-electron}
\end{equation}
With the nuclei clamped, $n_\kappa$ is a static Bragg sum, $\tilde R_{e\kappa} = 0$, and this is \cref{eq:4.4-eps} again. With moving nuclei the second term of $\tilde R_{eZ}$ is the electron density that follows a nucleus, and it is large at phonon energies, where the electrons of an atom move with their nucleus (\cref{sec:8.6}). The exact statement for this coupling is the electron-channel form of \cref{app:A.7}, $S_{ee} = (2q^2/e^2)\,L_e/(1 - \ee^{-\beta\omega})$ with $\varepsilon_e^{-1} = 1 + v_{\mathrm{C}}R_{ee}$, but $\varepsilon_e$ is not the dielectric function of any experiment. Its inverse has a pole at every acoustic phonon, weighted by the electrons that follow the nuclei, while the total $\varepsilon^{-1}$ at small $q$ has none, since a neutral cell that moves carries no charge.

\emph{Nucleon number.} A scalar that couples to protons and neutrons equally acts through $\sum_\kappa A_\kappa n_\kappa$, times its coupling per nucleon. Its vertex has no electronic part, and its screening runs entirely through $\tilde R_{NZ}$, the proper charge that moves when the nucleon density does. At phonon energies that charge is the Born charge of each sublattice; at electronic energies the nuclei do not move and the coupling is not screened at all. The Born charge never enters the vertex itself. It enters the pole: the proper response has its optical pole at the transverse frequency, the screened one at the longitudinal frequency, with the same f-sum weight (\cref{we:8-nucleon}).

\emph{Spin.} The spin density $\vect{s}_e$ of a target without spin--orbit coupling and without magnetic order is not screened. Its Hamiltonian commutes with the total spin, the charge density is invariant under spin rotations and $\vect{s}_e$ is a vector under them, and the equilibrium state is invariant too. The cross response $\tilde R_{sZ}$ then vanishes, and $R_{ss} = \tilde R_{ss}$. Spin--orbit coupling or magnetic order breaks the argument; the polarized gas of \cref{we:8-cross} has $S_{ns} \ne 0$ and screens its spin rows with its density. The nuclear spin density $\vect{s}_\kappa$ is unscreened for the same reason. Beyond independent particles a spin-flip pair still feels the attraction between its electron and hole, the direct kernel of \cref{sec:4.9}, but not the exchange kernel that carries the local fields: it has no charge to exchange.

The rows that are not densities follow the same rule. The longitudinal current is screened as the density. The continuity equation in the form of \cref{eq:8.1-transverse} makes $\hat{\vect{q}}\cdot\vect{j}_s$ act as $(\omega/q)\,n_s$ on the energy shell, and the Ward identity \cref{eq:4.4-ward} holds for proper and full responses alike. The transverse current and the kinetic-energy density $\mathcal{T}_e$ are screened where their proper correlation with the charge is nonzero; the spin current, a vector under spin rotations like $\vect{s}_e$, has none without spin--orbit coupling or magnetic order and is unscreened. The kinetic-energy density is a scalar under spin rotations and the point group and correlates with the charge already at $\vect{G} = 0$, so it is screened at the head, by a term of order $q^2$ at small $q$ in a crystal without a polar axis. In a cubic crystal, with $\vect{q}$ along a symmetry axis, $\tilde R_{\vect{j}_{\mathrm{T}}Z}$ vanishes at $\vect{G} = 0$ by symmetry, and the transverse current is screened through the local fields at $\vect{G} \ne 0$. In a crystal of lower symmetry, or off the symmetry axes, longitudinal and transverse parts mix already at $\vect{G} = 0$. Exchange of a transverse photon between currents is a separate and smaller effect at the momentum transfers of scattering: its propagator $1/(\omega^2 - q^2 - \Pi_{\mathrm{em}})$, with $\Pi_{\mathrm{em}}$ the photon self-energy of the medium (\cref{sec:8.4-mixing}), differs from $-1/q^2$ by $\Pi_{\mathrm{em}}/q^2 \sim \omega_p^2/q^2$, which is $2\times10^{-5}$ in silicon at $q_0$. Every row with $\tilde R_{aZ} \ne 0$ is screened through \cref{eq:8.4-screened}, and a rate that contracts \cref{eq:8.2-rate} with the proper response $\tilde R$ alone, or with a factor $1/|\varepsilon|^2$ in place of the cross terms, is wrong for every coupling except the charge.

\subsection{Exactness of the loss function}\label{sec:8.4-exact}

The loss function is measured for many materials, optically at $q \to 0$ and by electron energy loss at finite $q$, and it is the quantity that a computed response is checked against. The screened response fixes which couplings it describes exactly.
\begin{keyresult}[Exactness of the loss-function form]
Let a dark-matter particle couple to $\sum_sc_1^{(s)}n_s(\vect{q})$ in a crystal, a target whose long-wavelength longitudinal acoustic phonon moves it locally as a whole. Let nuclei of different atomic number have different masses, and set the nuclear form factors to one. Then
\begin{equation}
  S_{cc}(\vect{q},\omega) = \lambda_Z^2\,\frac{2q^2}{e^2}\,\frac{L(\vect{q},\omega)}{1 - \ee^{-\beta\omega}} \quad\text{at every } \vect{q} \text{ and every } \omega \neq 0 ,
  \label{eq:8.4-exact}
\end{equation}
with $S_{cc}$ the structure factor of the coupling and $L$ the loss function of the total dielectric matrix. For some constant $\lambda_Z$ it holds if and only if the coupling is proportional to the charge, $c_1^{(e)} = -\lambda_Z$ and $c_1^{(\kappa)} = \lambda_ZZ_\kappa$, or the nuclei are clamped and $c_1^{(e)} = \pm\lambda_Z$. A coupling to the electron density is exact only at a static lattice.
\end{keyresult}
\noindent Sufficiency is immediate. If the coupling differs from $\lambda_Zn_Z$ by an operator with no dynamics, a sum of static Bragg terms when the nuclei are clamped, the two have the same transitions and $S_{cc} = \lambda_Z^2S_{ZZ}$ at every $\omega \neq 0$, which is \cref{eq:8.4-exact} by the total-charge relation above.

The converse takes two limits in which electrons and nuclei come apart, and \cref{app:D.27} gives each step. At large momentum transfer, above the inverse size of every electron orbital and in the regime where a nucleus recoils freely from its site (the impulse regime of \cref{sec:8.5}), each particle recoils on its own. The structure factor of the coupling is then a sum of recoil peaks, at $q^2/2m_e$ for the electrons and $q^2/2M_\kappa$ for each nuclear species, with weights $(c_1^{(s)})^2\nd_s$, and that of the charge has the weights $\nd_e$ and $Z_\kappa^2\nd_\kappa$. Equality peak by peak requires $|c_1^{(e)}| = |\lambda_Z|$ and $|c_1^{(\kappa)}| = |\lambda_Z|Z_\kappa$ for every species: the magnitudes are those of the charge, but the signs are still free. The signs are fixed at the opposite end, $\vect{q} \to 0$ at the frequency of a longitudinal acoustic phonon. There the crystal moves as a whole, electrons included, and the amplitude to excite the phonon is proportional to the coupling carried by a unit volume, $\sum_sc_1^{(s)}\nd_s$. For the charge this sum is the net charge density, zero. Equality at small $q$ needs $c_1^{(e)}\nd_e + \sum_\kappa c_1^{(\kappa)}\nd_\kappa = 0$, and with the magnitudes already fixed this is possible only if every nuclear coefficient has the sign opposite to the electron coefficient. The coupling is then $\pm\lambda_Zn_Z$. With the nuclei clamped, the nuclear peaks and the acoustic phonon are absent, and only $|c_1^{(e)}| = |\lambda_Z|$ survives.

The theorem leaves no exact use of the loss function for an electrophilic mediator once the lattice moves, and none at all for a nucleon or isospin-violating coupling. Its practical content is a bound on where the error sits. For the electron density, the difference between $S_{ee}$ and $S_{ZZ}$ comes from the nuclear rows of \cref{eq:8.4-electron}. Above the phonon band, however, the nuclei reach the electronic final states only through the Migdal amplitudes, by which a moving nucleus excites the electrons (\cref{sec:8.6}), suppressed by powers of $\kappa_{\mathrm{BO}}$, and the loss function gives the electron-only rate to that accuracy, the statement of \cref{sec:4.4}. Inside the phonon band the two differ at leading order: $S_{ee}$ has an acoustic peak carrying the electrons that follow the nuclei, a fraction $Z_\kappa m_e/M_\kappa$ of its f-sum (\cref{eq:8.6-fsum-e}), while $S_{ZZ}$ has none at small $q$. \Cref{sec:10.2} uses the theorem for the electronic rate, and the relaxed-lattice correction to its sum-rule bound comes from this acoustic peak.

\Cref{tab:8-screening} collects the couplings with the screening each takes.

\begin{table}[tbp]
  \caption{External couplings of dark matter and the screening each takes, from the screened response \cref{eq:8.4-screened}. ``Exact $L$'' states whether the structure factor of the coupling equals $(2q^2/e^2)L/(1 - \ee^{-\beta\omega})$, with $L$ the loss function of the total dielectric matrix of electrons and nuclei, at every momentum transfer and energy (\cref{eq:8.4-exact}). SOC: spin--orbit coupling.}
  \label{tab:8-screening}
  \centering\small
  \renewcommand{\arraystretch}{1.25}
  \begin{tabular}{@{}>{\raggedright\arraybackslash}p{0.17\linewidth}>{\raggedright\arraybackslash}p{0.2\linewidth}>{\raggedright\arraybackslash}p{0.35\linewidth}>{\raggedright\arraybackslash}p{0.19\linewidth}@{}}
    \toprule
    Coupling & Example & Screened by & Exact $L$ \\
    \midrule
    Total charge $n_Z$ & dark photon, heavy or light; millicharge & $\varepsilon^{-1}$ of electrons and nuclei; $v_{\mathrm{C}}R_{ZZ} = \varepsilon^{-1} - 1$ & yes, any lattice \\
    Electron density $n_e$ & electrophilic scalar or vector mediator & $\tilde R_{eZ}\varepsilon^{-1}v_{\mathrm{C}}\tilde R_{Ze}$: electrons, and the nuclei they follow & clamped nuclei, $\omega \ne 0$; above the phonon band to Migdal order \\
    Nucleon number $\sum_\kappa A_\kappa n_\kappa$ & hadrophilic scalar & ionic polarization through the Born charges; it moves the poles (TO to LO), not the vertex & no \\
    Other density mixtures & isospin-violating couplings & $\tilde R_{cZ}\varepsilon^{-1}v_{\mathrm{C}}\tilde R_{Zc}$ & only if proportional to $n_Z$ \\
    Spin density $\vect{s}_e$, $\vect{s}_\kappa$ & $\mathcal{O}_4$, $\mathcal{O}_6$, $\mathcal{O}_{10}$ & none without SOC or magnetic order; the triplet exchange--correlation kernel remains & no \\
    Longitudinal current $\hat{\vect{q}}\cdot\vect{j}$ & $\mathcal{O}_8$, $\mathcal{O}_2$ (parts along $\vect{q}$) & as the density, $(\omega/q)$ times its row & as the density \\
    Transverse current $\vect{j}_{\mathrm{T}}$ & $\mathcal{O}_5$, $\mathcal{O}_8$ (part across $\vect{q}$) & only through local fields ($\vect{G} \ne 0$) on symmetry axes; at $\vect{G} = 0$ only through longitudinal--transverse mixing & no \\
    Spin current $\mathcal{J}^{\alpha\beta}_e$ & $\mathcal{O}_3$, $\mathcal{O}_7$, $\mathcal{O}_{12}$--$\mathcal{O}_{15}$; ALP absorption & none without SOC or magnetic order, as the spin density & no \\
    Kinetic-energy density $\mathcal{T}_e$ & $\mathcal{O}_2$ (part in $(\vect{v}^\perp)^2$); scalar absorption & head and local fields; head term $O(q^2)$ at small $q$ in crystals without a polar axis & no \\
    \bottomrule
  \end{tabular}
\end{table}

\begin{workedexample}{The nucleon-number coupling of a heavy scalar in GaAs}\label{we:8-nucleon}
Take a heavy scalar that couples to protons and neutrons with the same coefficient, $c_1^{(p)} = c_1^{(n)}$, and to no electron. By \cref{eq:8.1-nuclear} its row is $c_1^{(p)}\sum_\kappa A_\kappa n_\kappa(\vect{q})$, with $F_N = 1$ at the momenta of phonon scattering. We use the measured GaAs numbers of \cref{we:5-lst}: $\omega_{\mathrm{T}} = \SI{33.62}{meV}$, $\omega_{\mathrm{L}} = \SI{36.35}{meV}$, $\varepsilon_\infty = 10.88$, $\Omega = \SI{44.80}{\angstrom^3}$, $\mu_{\mathrm{red}} = \SI{36.11}{u}$ and $Z^* = 2.11$. Take $\vect{q}$ small against the zone and $\omega$ near the optical phonons, where the macroscopic field is the whole of the long-range Coulomb line and \cref{eq:8.4-rank1} applies.

\emph{Vertices.} In an optical mode at small $\vect{q}$ the two sublattices move against each other with the center of mass at rest, $\vect{u}_{\mathrm{Ga}} = (\mu_{\mathrm{red}}/M_{\mathrm{Ga}})\vect{u}_{\mathrm{rel}}$ and $\vect{u}_{\mathrm{As}} = -(\mu_{\mathrm{red}}/M_{\mathrm{As}})\vect{u}_{\mathrm{rel}}$, with $\vect{u}_{\mathrm{rel}} = \vect{u}_{\mathrm{Ga}} - \vect{u}_{\mathrm{As}}$ the relative coordinate of \cref{sec:5.5}. Expanding $n_\kappa(\vect{q}) = \sum_{\vect{R}}\ee^{\ii\vect{q}\cdot(\vect{R} + \vect{\tau}_\kappa + \vect{u}_{\vect{R}\kappa})}$ to first order in the displacements, the nucleon row creates the mode through $\ii\vect{q}\cdot\vect{u}_{\mathrm{rel}}$ times an amplitude per cell, which at leading order in $\vect{q}$ is
\begin{equation*}
  a_N = \mu_{\mathrm{red}}\Bigl(\frac{A_{\mathrm{Ga}}}{M_{\mathrm{Ga}}} - \frac{A_{\mathrm{As}}}{M_{\mathrm{As}}}\Bigr) .
\end{equation*}
At finite $\vect{q}$ the site phases $\ee^{\ii\vect{q}\cdot\vect{\tau}_\kappa}$ of $n_\kappa$, combined with the zone-center displacements, give a vertex that changes when another arsenic atom is assigned to the cell. Only with the eigenvector of the mode at $\vect{q}$ is the vertex independent of that choice, $a_N(\vect{q}) = \mu_{\mathrm{red}}^{1/2}\sum_\kappa A_\kappa\ee^{\ii\vect{q}\cdot\vect{\tau}_\kappa}\,\hat{\vect{q}}\cdot\vect{e}^*_{\kappa\mathrm{L}}(\vect{q})/M_\kappa^{1/2}$, with the conventions of \cref{eq:11.1-S1}.
The charge row creates it through the Born charges, $\sum_\kappa Z^*_\kappa\,\ii\vect{q}\cdot\vect{u}_\kappa = Z^*\,\ii\vect{q}\cdot\vect{u}_{\mathrm{rel}}$ at leading order in $\vect{q}$, since $\mu_{\mathrm{red}}(1/M_{\mathrm{Ga}} + 1/M_{\mathrm{As}}) = 1$. The Born charge is the right vertex for the proper response: \cref{sec:5.5} defines it at zero macroscopic field, which is what removing the Coulomb line at $\vect{q}$ imposes. For the same reason the relative coordinate of the proper response oscillates at $\omega_{\mathrm{T}}$, for longitudinal as well as transverse motion, with the response $1/[\mu_{\mathrm{red}}(\omega^2 - \omega_{\mathrm{T}}^2)]$ per cell. With $N/\Vol = 1/\Omega$,
\begin{equation*}
  \tilde R_{NN} = \frac{q^2|a_N|^2}{\mu_{\mathrm{red}}\Omega\,(\omega^2 - \omega_{\mathrm{T}}^2)} , \qquad
  \tilde R_{NZ} = \frac{q^2a_NZ^*}{\mu_{\mathrm{red}}\Omega\,(\omega^2 - \omega_{\mathrm{T}}^2)} , \qquad
  \varepsilon = \varepsilon_\infty + \frac{e^2Z^{*2}}{\Omega\,\mu_{\mathrm{red}}(\omega_{\mathrm{T}}^2 - \omega^2)} ,
\end{equation*}
where $\varepsilon$ is \cref{eq:5.5-eps}: the electrons give $\varepsilon_\infty$, and the ionic part of $\tilde R_{ZZ}$ the rest.

\emph{Screening.} Insert these into \cref{eq:8.4-rank1}. By \cref{eq:5.5-eps,eq:5.5-LOTO}, $\varepsilon = \varepsilon_\infty(\omega_{\mathrm{L}}^2 - \omega^2)/(\omega_{\mathrm{T}}^2 - \omega^2)$ and $e^2Z^{*2}/(\Omega\mu_{\mathrm{red}}\varepsilon_\infty) = \omega_{\mathrm{L}}^2 - \omega_{\mathrm{T}}^2$, and the screening term is
\begin{equation*}
  v_{\mathrm{C}}\frac{\tilde R_{NZ}\tilde R_{ZN}}{\varepsilon} = \frac{q^2|a_N|^2}{\mu_{\mathrm{red}}\Omega}\,\frac{\omega_{\mathrm{L}}^2 - \omega_{\mathrm{T}}^2}{(\omega^2 - \omega_{\mathrm{T}}^2)(\omega^2 - \omega_{\mathrm{L}}^2)} ,
\end{equation*}
and added to $\tilde R_{NN}$ it cancels the transverse pole:
\begin{equation*}
  R_{NN}(\vect{q},\omega) = \frac{q^2|a_N|^2}{\mu_{\mathrm{red}}\Omega\,(\omega^2 - \omega_{\mathrm{L}}^2)} , \qquad
  S_{NN}(\vect{q},\omega) = \frac{\pi q^2|a_N|^2}{\mu_{\mathrm{red}}\Omega\,\omega_{\mathrm{L}}}\,\delta(\omega - \omega_{\mathrm{L}}) ,
\end{equation*}
at $T = 0$ and $\omega > 0$. The screened response is the proper one with its pole moved from $\omega_{\mathrm{T}}$ to $\omega_{\mathrm{L}}$, and its f-sum weight $\int\dd\omega\,\omega S_{NN}/2\pi$ is unchanged; the area of the peak, $\int\dd\omega\,S_{NN}/2\pi$, falls by $\omega_{\mathrm{T}}/\omega_{\mathrm{L}} = 0.925$. The Born charges and $\varepsilon_\infty$ enter through $\omega_{\mathrm{L}}$ and nowhere else; the vertex is $a_N$, a property of the nuclei and their masses. The rule $1/|\varepsilon|^2$ gets this wrong in an instructive way. The proper structure factor has its weight at $\omega_{\mathrm{T}}$, where $\varepsilon$ has a pole and $1/|\varepsilon|^2 = 0$, so the naive rate would vanish; the true rate sits at $\omega_{\mathrm{L}}$, where $1/|\varepsilon|^2$ is infinite. The cross term of \cref{eq:8.4-rank1}, which the naive rule drops, resolves both.

\emph{The size of the vertex.} At $\vect{q} \to 0$ the phases drop out and $a_N \to \mu_{\mathrm{red}}(A_{\mathrm{Ga}}/M_{\mathrm{Ga}} - A_{\mathrm{As}}/M_{\mathrm{As}})$, which would vanish if atomic masses were proportional to nucleon numbers. They nearly are. With the atomic masses of the isotopes~\cite{wang2021ame2020} and the natural abundances of gallium, \SI{60.108}{\percent} $^{69}$Ga and \SI{39.892}{\percent} $^{71}$Ga~\cite{meija2016isotopic}, the mean nucleon number per mass is $1.001072$ per u for gallium and $1.001046$ per u for arsenic, and $a_N = 36.11 \times 2.59\times10^{-5} = 9.4\times10^{-4}$. A nucleon-number coupling barely moves the center of nucleon number when the center of mass stays put. At finite $\vect{q}$ the center of mass no longer stays put. In zincblende a shear strain displaces the two sublattices against each other, the internal strain behind the piezoelectric charge of \cref{sec:6.2}, and by the same coupling an optical wave drags the center of mass along wherever a longitudinal wave contains shear. The nucleon row then sees the mass density of the cell, $A_\kappa/M_\kappa \approx 1$ per u, and a coupling to the mass density alone gives nearly the same vertex. Along $[111]$, the direction of the Ga--As bond, a longitudinal wave contains shear, and the drag grows linearly in $q$ and overtakes the mass term near $q \approx \SI{0.05}{eV}$. Scattering reaches $\omega_{\mathrm{L}}$ only for $q \ge \omega_{\mathrm{L}}/v_{\mathrm{max}} = \SI{13.7}{eV}$, where, with the phonons of \cref{fig:5-gaas} at that $\vect{q}$, $|a_N| = 0.24$, 250 times the long-wavelength value; at \SI{100}{eV}, $|a_N| = 1.74$. Along $[100]$ a longitudinal wave contains no shear, the drag is absent at first order in $q$, and $|a_N| = 2.1\times10^{-3}$ at \SI{13.7}{eV}, close to its value at $q \to 0$. Compared with the longitudinal acoustic phonon, whose vertex $\sum_\kappa A_\kappa = 144.8$ per cell is not screened at leading order because $\sum_\kappa Z^*_\kappa = 0$, the energy-weighted optical share, $|a_N|^2M_{\mathrm{cell}}/\mu_{\mathrm{red}}$ divided by the squared acoustic vertex, grows as $q^2$; along $[111]$ it is $1.1\times10^{-5}$ at \SI{13.7}{eV} and $5.8\times10^{-4}$ at \SI{100}{eV}. This is the suppression of optical phonons by destructive interference within the cell that Griffin and collaborators found for couplings proportional to the mass number~\cite[Eqs.~(30) and (33)]{griffin2018directional}. For a dark photon the same mode couples through $Z^*/\varepsilon_\infty = 0.19$ at any small $q$, and it is the polar phonon of \cref{sec:11.2}.
\end{workedexample}

\subsection{Coulomb-line classes}\label{sec:8.4-lines}

Below the direct gap of silicon, dark matter can excite an electron--hole pair only together with a phonon (\cref{sec:6.5}). The amplitude contains three objects that the medium screens. The dark-matter vertex sits at $(\vect{q},\omega)$. The electron--phonon vertex sits at the phonon wavevector $\vect{Q}$, and density-functional perturbation theory returns it already screened (\cref{sec:6.2}). The pair that leaves has momentum $\vect{q} + \vect{Q}$ and energy $\omega \mp \omega_{\vect{Q}\nu}$, upper sign for emission. A calculation that divides the dark-matter vertex by $\varepsilon(\vect{q},\omega)$ and uses the computed vertex $g$ might count one Coulomb line twice, or miss one. The question is which Coulomb lines each factor contains, and the diagrams of \cref{sec:4.8} answer it.

Work at the random-phase level, where every Coulomb line joins two different fermion loops and nothing runs inside a loop, at second order in the electron--phonon vertex $g$ in the rate, with one phonon in the final state. Draw a diagram of the response as loops joined by Coulomb lines, with one phonon line between two ion vertices and the two dark-matter vertices on loops.
\begin{keyresult}[Coulomb-line classes]
In such a diagram every Coulomb line belongs to exactly one of three classes.
\begin{itemize}
  \item \emph{External}, at $(\vect{q} + \vect{G}, \omega)$: a line whose removal cuts the diagram in two. These lines form the chains that \cref{eq:8.4-screened} resums into $(\varepsilon^{-1})_{\vect{G}0}(\vect{q},\omega)$ on the dark-matter vertex, local fields included.
  \item \emph{Vertex}, at $\vect{Q} + \vect{G}$ and frequencies of order $\omega_{\vect{Q}\nu}$: lines between an ion vertex and the first loop that carries a dark-matter vertex. They make the screened electron--phonon vertex $g$.
  \item \emph{Pair}, at $(\vect{q} + \vect{Q} + \vect{G}, \omega \mp \omega_{\vect{Q}\nu})$: lines between the two loops that carry the dark-matter vertices. They resum into the screened interaction $W^{\mathrm{scr}}(\vect{q} + \vect{Q}, \omega \mp \omega_{\vect{Q}\nu})$ of \cref{sec:4.8}, acting on the final pair.
\end{itemize}
A phonon line that itself cuts the diagram in two is not a Coulomb line. It is the ionic part of the total dielectric matrix, included by using the $\varepsilon$ of \cref{eq:8.4-screened} in the external and pair chains.
\end{keyresult}
\noindent The proof is by the topology of the diagram (\cref{app:D.27}). With the phonon line removed, the loops and Coulomb lines form a tree, since at the random-phase level every Coulomb line is a bridge; a closed cycle of Coulomb lines would need a line inside a loop. Putting the phonon line back closes exactly one cycle. Lines off the cycle remain bridges, and momentum conservation sets them to $\vect{q} + \vect{G}$; these are the external lines. On the cycle, the segments next to the ion vertices carry $\vect{Q} + \vect{G}$, and the segment between the two loops with dark-matter vertices carries $\vect{q} + \vect{Q} + \vect{G}$. When both dark-matter vertices sit on one loop, the pair segment is absent. The class of a line follows from its position, not from its momentum label, and the classes stay distinct where the labels coincide. At $\vect{q} \to 0$ the pair segment carries the momentum $\vect{Q}$ of the vertex segments but an energy of order eV against theirs of order $\omega_{\vect{Q}\nu}$, and at $\vect{Q} \to 0$ the pair segment lies on the cycle while the external lines do not.

Three consequences follow. First, the screened vertex of density-functional perturbation theory, an external $\varepsilon^{-1}(\vect{q},\omega)$ and the pair interaction $W^{\mathrm{scr}}$ together count every random-phase diagram once. An external $1/\varepsilon(\vect{Q})$ on top of the computed $g$ would count the vertex lines twice; leaving out $W^{\mathrm{scr}}$ is an omission, not a cure for a double count. Second, the vertex lines are evaluated by density-functional perturbation theory at zero frequency, while they carry frequencies of order $\omega_{\vect{Q}\nu}$. Below the gap, time reversal makes $\varepsilon(\vect{Q},\omega')$ an even and analytic function of the frequency $\omega'$ for $|\omega'| < \Eg$ (\cref{app:D.27}). The error of the static vertex is then of order $(\omega_{\vect{Q}\nu}/\Eg)^2$, about $3\times10^{-3}$ in silicon, the number of \cref{we:8-label}; in a magnetic insulator, without time reversal, it can be of order $\omega_{\vect{Q}\nu}/\Eg$. Third, the ionic insertion is small at electronic energies. Far above the optical phonons, \cref{eq:5.5-eps} gives an ionic part $\varepsilon - \varepsilon_\infty \approx -(\varepsilon_0 - \varepsilon_\infty)\,\omega_{\mathrm{T}}^2/\omega^2$, which is $-0.9\times10^{-3}$ in GaAs at \SI{1.5}{eV}, $10^{-4}$ of $\varepsilon_\infty$; in silicon, whose Born charges vanish, it is smaller still.

The pair lines are the new element. They screen the final pair at its own momentum and energy: once the phonon has taken $\vect{Q}$, the charge that the dark-matter vertex induced oscillates at $\vect{q} + \vect{Q}$, and the medium screens it there, not at $\vect{q}$. The final pair is on shell, the screened interaction at its momentum and energy is complex, and it acts between the two halves of the amplitude, which \cref{sec:13.4} writes out with its size in silicon. Lines inside a loop, the electron--hole ladder of \cref{sec:4.9}, lie outside the classification; an adiabatic exchange--correlation kernel adds a source term of its own. A frozen-phonon supercell calculation with local fields in the supercell basis contains all three classes at once, since every $\vect{Q} + \vect{G}$ is then a reciprocal vector of the supercell.

\subsection{In-medium mixing}\label{sec:8.4-mixing}

At \SI{3}{eV} silicon has $\varepsilon_2 = 2.4$ and a loss function $L = 3.1\times10^{-3}$ (\cref{sec:4.4}), so $|\varepsilon|^2 = \varepsilon_2/L = 761$. A dark photon absorbed there couples as if its squared mixing were $\zeta^2/761$ rather than $\zeta^2$. The factor is the in-medium mixing $\zeta_{\mathrm{eff}}$ of \cref{sec:1.4}, and it comes from the photon-reducible part of the response. A dark photon couples to the electromagnetic current with strength $\zeta e$ (\cref{sec:1.4}). The current it induces radiates an ordinary photon, which the medium absorbs again, and the chain repeats. Split the current response as the density response was split: a part $\tilde R$ that is proper with respect to one photon line, longitudinal or transverse, and the photon-reducible chains. For absorption at rest, \cref{app:D.25} summed the chain and found that it turns $\zeta^2$ into
\begin{equation}
  \zeta_{\mathrm{eff}}^2 = \frac{\zeta^2m_\phi^4}{|m_\phi^2 - \Pi_{\mathrm{em}}|^2} ,
  \label{eq:8.4-zetaeff}
\end{equation}
with $\Pi_{\mathrm{em}}$ the photon self-energy, the proper current response times $e^2$, diamagnetic term included. In the language of Mitridate and collaborators, the dark-matter self-energy acquires a mixing term $\Pi_{\phi A}\Pi_{A\phi}/(m_\phi^2 - \Pi_{AA})$ beside its proper part~\cite[Eqs.~(20)--(22)]{mitridate2021dark}; that term is the photon-reducible response.

Both polarizations follow from one relation. The transverse self-energy is $\Pi_{\mathrm{T}} = \omega^2(1 - \varepsilon^{\mathrm{T}})$ (\cref{sec:4.4}), and the longitudinal one is $\Pi_{\mathrm{L}} = (\omega^2 - k^2)(1 - \varepsilon)$~\cite[Eq.~(34)]{mitridate2021dark}, with $k$ the momentum of the dark photon and $\varepsilon$ the longitudinal function of \cref{eq:4.4-eps}. On shell, $\omega^2 - k^2 = m_\phi^2$, and $m_\phi^2 - \Pi_{\mathrm{L}} = m_\phi^2\varepsilon$: a longitudinal dark photon mixes with $\zeta_{\mathrm{eff}} = \zeta/\varepsilon$ exactly. A transverse one does so at $k \to 0$, where $\varepsilon^{\mathrm{T}} = \varepsilon$ by \cref{eq:4.4-ET} and $\omega = m_\phi$. The net rate per quantum, \cref{eq:8.2-abs} with the full response, is then
\begin{equation*}
  \varGamma_{\mathrm{net}} = -\frac{\zeta^2m_\phi^2}{\omega}\,\Imag\frac{\Pi_{\mathrm{em}}}{m_\phi^2 - \Pi_{\mathrm{em}}} = -\zeta^2m_\phi\,\Imag\Bigl(\frac{1}{\varepsilon} - 1\Bigr) = \zeta^2m_\phi\,L(m_\phi) ,
\end{equation*}
which is $\zeta_{\mathrm{eff}}^2\sigma_1$ with $\sigma_1 = \omega\varepsilon_2$~\cite[Eq.~(40)]{mitridate2021dark}. A dark photon at rest measures the loss function of the target, not its absorption. This is why its absorption in GaAs peaks at the longitudinal phonon (\cref{sec:5.5}).

Scattering through a dark-photon mediator is the same construction at $\omega \ll q$. The interaction of \cref{we:8-darkphoton} uses the potential of the dark matter in vacuum, $\tilde V$, and the full response $R_{ZZ}$, which carries the chain through $\varepsilon^{-1}$. We can instead put the chain into the mediator: the dark photon then mixes with the longitudinal photon of the medium, its propagator acquires $1/\varepsilon(\vect{q},\omega)$, and it acts on the proper response. The two bases agree~\cite[Eq.~(8) and the text below it]{knapen2021dark}: with $\varepsilon = 1 - v_{\mathrm{C}}\tilde R_{ZZ}$, $\tilde R_{ZZ}/\varepsilon = \tilde R_{ZZ}\varepsilon^*/|\varepsilon|^2$, and $\Imag(\tilde R_{ZZ}\varepsilon^*) = \Imag\tilde R_{ZZ}$ because $v_{\mathrm{C}}|\tilde R_{ZZ}|^2$ is real, so that $\Imag R_{ZZ} = \Imag\tilde R_{ZZ}/|\varepsilon|^2$. In scattering, the in-medium mixing is the factor $1/\varepsilon(\vect{q},\omega)$, and it is the screening of the total-charge row of \cref{tab:8-screening}.

The bookkeeping rule is the one of the Coulomb lines. Either the vacuum mixing $\zeta$ multiplies the full response, or $\zeta_{\mathrm{eff}}$ multiplies the proper one. The loss function already contains the screening, and a factor $\zeta_{\mathrm{eff}}^2/\zeta^2$ applied to a loss-function rate counts it twice. The optical conductivity $\sigma_1$ is a proper response, and \cref{eq:1.4-dp-abs} multiplies it by $\zeta_{\mathrm{eff}}^2$ once. Other bosons mix in the same way through their own rows. The scalar of \cref{eq:8.1-scalar} couples to $n_e$ and $\mathcal{T}_e$, both of which correlate with the charge, and its absorption carries a photon-reducible part through $\tilde R_{n_eZ}$ and $\tilde R_{\mathcal{T}_eZ}$. An axion-like particle in a static field $\vect{B}_0$ acts as a source current along $\vect{B}_0$ and mixes as a dark photon does; the spin rows of the axioelectric coupling do not mix at all without spin--orbit coupling. \Cref{sec:10.7} works out these rates.
